\documentclass[11pt]{article}

\usepackage[margin=1in]{geometry}
\usepackage{amsmath}
\usepackage{amssymb}
\usepackage[T1]{fontenc}
\usepackage[utf8]{inputenc}

\title{Radiation Reaction Model Supplement: Native-Units Mapping}
\author{LW Integrator Notes}
\date{\today}

\begin{document}

\maketitle

\section{Purpose}

This note restates the candidate Medina radiation-reaction model in the
solver's native unit system and contrasts it with the currently implemented
\texttt{power\_matched\_damping} mode.

\section{Native Units}

The maintained solver uses the historical native system:

\begin{align}
x &: \mathrm{mm}, &
t &: \mathrm{ns}, &
v &: \mathrm{mm/ns},
\\
m &: \mathrm{amu}, &
p &: \mathrm{amu\,mm/ns},
\\
F &: \mathrm{amu\,mm/ns^2}, &
E_{\mathrm{mech}} &: \mathrm{amu\,mm^2/ns^2}.
\end{align}

The native elementary charge used by the solver is derived from the
statcoulomb conversion:
\begin{equation}
e_{\mathrm{native}}
=e_{\mathrm{statC}}C_{\mathrm{statC}\rightarrow\mathrm{native}}
\simeq 1.1787088598 \times 10^{-5}.
\end{equation}

A native electric field is a force-per-charge quantity, so
\begin{equation}
qE \sim \mathrm{amu\,mm/ns^2}.
\end{equation}

\section{Time Variable}

The solver step is proper time:
\begin{equation}
h = d\tau.
\end{equation}

Coordinate time advances as
\begin{equation}
dt = \gamma h.
\end{equation}

Mechanical momentum is
\begin{equation}
\mathbf{p} = \gamma m \mathbf{v}
            = \gamma m c \boldsymbol{\beta}.
\end{equation}

The current prescribed external-field impulse is
\begin{equation}
\Delta \mathbf{p}_{\mathrm{ext}}
= h q \gamma \left(\mathbf{E} + \boldsymbol{\beta} \times \mathbf{B}\right).
\end{equation}

Since \(dt = \gamma h\), this equals
\begin{equation}
\Delta \mathbf{p}_{\mathrm{ext}}
= dt\,\mathbf{F}_{\mathrm{ext}},
\end{equation}
where
\begin{equation}
\mathbf{F}_{\mathrm{ext}}
= q\left(\mathbf{E} + \boldsymbol{\beta} \times \mathbf{B}\right).
\end{equation}

\section{Medina Force}

The Medina radiation-reaction force from the paper is
\begin{equation}
\mathbf{F}_{\mathrm{RAD}}
=
\frac{2}{3}\frac{q^2}{mc^3}
\left[
\frac{d}{dt}\left(\gamma\mathbf{F}_{\mathrm{ext}}\right)
-
\frac{\gamma^3}{c^2}
\left(\mathbf{F}_{\mathrm{ext}}\cdot \mathbf{a}\right)
\mathbf{v}
\right].
\end{equation}

Define
\begin{equation}
\tau_0 = \frac{2}{3}\frac{q^2}{mc^3}.
\end{equation}

Then
\begin{equation}
\mathbf{F}_{\mathrm{RAD}}
=
\tau_0
\left[
\frac{d}{dt}\left(\gamma\mathbf{F}_{\mathrm{ext}}\right)
-
\frac{\gamma^3}{c^2}
\left(\mathbf{F}_{\mathrm{ext}}\cdot \mathbf{a}\right)
\mathbf{v}
\right],
\end{equation}
with
\begin{align}
\mathbf{v} &= c\boldsymbol{\beta}, \\
\mathbf{a} &= \frac{d\mathbf{v}}{dt}
            = c\frac{d\boldsymbol{\beta}}{dt}.
\end{align}

The complete derivative is
\begin{equation}
\frac{d}{dt}\left(\gamma\mathbf{F}_{\mathrm{ext}}\right)
=
\gamma\frac{d\mathbf{F}_{\mathrm{ext}}}{dt}
+
\frac{d\gamma}{dt}\mathbf{F}_{\mathrm{ext}}.
\end{equation}
The first implementation accidentally retained only the second term.  That
specialization is valid for a constant external force, but not for a Coulomb
flyby or another spatially varying field.

The Medina impulse over one solver step would be
\begin{equation}
\Delta \mathbf{p}_{\mathrm{RAD}}
=
\mathbf{F}_{\mathrm{RAD}}\,dt
=
\mathbf{F}_{\mathrm{RAD}}\gamma h.
\end{equation}

Apply it to mechanical momentum:
\begin{equation}
\mathbf{p}_{\mathrm{new}}
=
\mathbf{p}_{\mathrm{nonrad}}
+
\Delta \mathbf{p}_{\mathrm{RAD}}.
\end{equation}

Then recompute gamma:
\begin{equation}
\gamma_{\mathrm{new}}
=
\sqrt{
1 +
\frac{|\mathbf{p}_{\mathrm{new}}|^2}{(mc)^2}
}.
\end{equation}

Then recompose canonical momentum:
\begin{align}
P_{i,\mathrm{new}}
&=
p_{i,\mathrm{new}}
+
m A_{i,\mathrm{solver}},
\\
P_{t,\mathrm{new}}
&=
\gamma_{\mathrm{new}}mc
+
q\Phi.
\end{align}

\section{Difference From Current Radiation-Reaction Mode}

The current \texttt{power\_matched\_damping} mode is not Medina/LAD. It is an
energy-bookkeeping approximation.

It computes radiated energy:
\begin{equation}
E_{\mathrm{rad}}
=
P_{\mathrm{Lienard}}\,dt.
\end{equation}

Then it removes that much kinetic energy by scaling the mechanical momentum
magnitude:
\begin{equation}
\mathbf{p}_{\mathrm{new}}
=
\mathbf{p}_{\mathrm{old}}
\frac{|\mathbf{p}_{\mathrm{target}}|}{|\mathbf{p}_{\mathrm{old}}|}.
\end{equation}

So it preserves direction:
\begin{equation}
\hat{\mathbf{p}}_{\mathrm{new}}
=
\hat{\mathbf{p}}_{\mathrm{old}}.
\end{equation}

Medina instead computes a vector force from
\begin{equation}
\mathbf{F}_{\mathrm{ext}}, \qquad
\frac{d\gamma}{dt}, \qquad
\mathbf{a}, \qquad
\mathbf{v}.
\end{equation}

Thus, Medina can change both the magnitude and direction of mechanical
momentum.

For the current experimental implementation, accepted interval-midpoint force
samples supply a first-order backward estimate of
\(d(\gamma\mathbf{F}_{\mathrm{ext}})/dt\).  The first unprimed sample records
far radiation but applies no incomplete reaction impulse.
The quantities \(d\boldsymbol{\beta}/dt\) and \(d\gamma/dt\) are derived from
the estimated mechanical force:
\begin{align}
\frac{d\boldsymbol{\beta}}{dt}
&=
\frac{\mathbf{F}_{\mathrm{ext}}
- \boldsymbol{\beta}(\boldsymbol{\beta}\cdot\mathbf{F}_{\mathrm{ext}})}
{\gamma m c},
\\
\frac{d\gamma}{dt}
&=
\frac{\boldsymbol{\beta}\cdot\mathbf{F}_{\mathrm{ext}}}{m c}.
\end{align}
This keeps the one-dimensional longitudinal terms self-consistent. For a
constant ideal collinear force, the two Medina bracket terms cancel to
numerical precision. In a transverse bend, \(d\gamma/dt=0\) while
\(\mathbf{F}_{\mathrm{ext}}\cdot\mathbf{a}\neq 0\), producing a recoil mostly
opposite the particle velocity.

\section{Practical Summary}

\begin{description}
\item[Current \texttt{power\_matched\_damping}]
Computes radiated energy, subtracts that energy from the mechanical momentum
magnitude, preserves momentum direction, and is best treated as bounded energy
bookkeeping.

\item[Medina/LAD candidate]
Computes the complete force derivative from accepted force samples, applies a
radiation-reaction impulse to mechanical momentum, may change both direction
and magnitude, and exposes derivative-readiness, cap, signed-work, radiated
energy, and cross-field-energy diagnostics.  It remains experimental until
active recoil converges for retarded source forces.
\end{description}

The first Medina implementation should use rest mass in \(\tau_0\) and
explicitly state that the paper's dressed-mass caveat is not yet modeled. It
has started with prescribed-field benchmarks before being applied to retarded
image-charge collision cases.

\end{document}
